\clearpage
\SourceLeaf{374}%
\begin{center}{\large [THE SURFACE OF THE ELLIPSOID.]\par}\end{center}
\begin{center}[1.]\end{center}

If each point of a surface is represented by the point on a sphere of radius 1 whose radius vector is parallel to the normal at the first point, and if an element $d\sigma$ of the sphere’s surface corresponds to an element $ds$ of the surface, then
\[
\frac{d\sigma}{ds}=\frac{\dfrac{ddz}{dx^2}\dfrac{ddz}{dy^2}-\left(\dfrac{ddz}{dx\,dy}\right)^2}{\left(1+\left(\dfrac{dz}{dx}\right)^2+\left(\dfrac{dz}{dy}\right)^2\right)^2}=\frac1{rr'},
\]
where $r$ and $r'$ are the two extreme radii of curvature.
\begin{center}[2.]\end{center}

If $d\sigma$ denotes the region on the celestial sphere that corresponds to the element $ds$ of the surface of the ellipsoid
\[
\left[\frac{xx}{aa}+\frac{yy}{bb}+\frac{zz}{cc}=1\right]
\]
then
\[
ds=aabbcc\left(\frac{xx}{a^4}+\frac{yy}{b^4}+\frac{zz}{c^4}\right)^2d\sigma.
\]
If $\xi,\eta,\zeta$ are the coordinates of the point on the celestial sphere, then
